% Figure: lifting padeye geometry for the end-to-end case study. A plate
% with a central pin hole, a load P pulling up through the pin, the main
% plate body below carrying tension to the weld at the base. Annotated
% with the net-section plane, the hole diameter, the cheek dimension, and
% the worst-stress points at the hole edge.
\begin{figure}[htbp]
\centering
\begin{tikzpicture}[scale=1.0]
% main plate outline (rounded top, rectangular base)
\draw[engblack, very thick]
  (-2.2,0) -- (-2.2,2.6) arc[start angle=180, end angle=0, radius=2.2]
  -- (2.2,0) -- (-2.2,0) -- cycle;
% base weld line
\draw[engred, very thick] (-2.2,0) -- (2.2,0);
\node[engred, font=\scriptsize, anchor=north] at (0,-0.1) {base weld to deck};
% pin hole, centred at (0, 2.6), diameter such that radius = 0.9
\draw[engblue, very thick, fill=white] (0,2.6) circle[radius=0.9];
\filldraw[enggray] (0,2.6) circle[radius=1.2pt];
% applied load arrow through the pin
\draw[engorange, very thick, -{Latex}] (0,2.6) -- (0,5.2);
\node[engorange, font=\small, anchor=west] at (0.15,4.7) {$P = 120\ \text{kN}$};
% hole diameter annotation
\draw[enggray, {Latex}-{Latex}] (-0.9,2.6) -- (0.9,2.6);
\node[enggray, font=\scriptsize, anchor=north] at (0,2.5) {$d_h = 35\ \text{mm}$};
% net-section plane through the hole (perpendicular to load)
\draw[englime!60!black, thick, dashed] (-2.6,2.6) -- (-0.9,2.6);
\draw[englime!60!black, thick, dashed] (0.9,2.6) -- (2.6,2.6);
\node[englime!60!black, font=\scriptsize, anchor=east] at (-2.65,2.6) {net section};
% worst-stress points at hole edge (on the load-perpendicular diameter)
\filldraw[engred] (-0.9,2.6) circle[radius=2pt];
\filldraw[engred] (0.9,2.6) circle[radius=2pt];
\node[engred, font=\scriptsize, anchor=south] at (-1.5,2.75) {peak $\sigma$};
% plate width annotation at net section
\draw[enggray, {Latex}-{Latex}] (-2.4,1.0) -- (2.4,1.0);
\node[enggray, font=\scriptsize, anchor=south] at (0,1.0) {$w = 110\ \text{mm}$};
% plate thickness note
\node[engblack, font=\scriptsize, anchor=west] at (2.4,0.4) {plate $t = 20\ \text{mm}$};
\end{tikzpicture}
\caption[Lifting padeye geometry]{The lifting padeye of the end-to-end
case study. A steel plate of thickness $t = 20\ \text{mm}$ and width
$w = 110\ \text{mm}$ at the net section carries a lift load $P = 120\
\text{kN}$ through a pin in a hole of diameter $d_h = 35\ \text{mm}$.
The load reaches the supporting structure through a base weld. The two
worst-stress points sit at the hole edge on the diameter perpendicular
to the load; the net section (dashed) is the reduced cross-section the
tension must cross. The forensic question is whether the peak stress at
those points survives repeated lifts.}
\label{fig:vol03ch03:padeye-geometry}
\end{figure}
